\documentclass{article}
\usepackage{amsmath, amssymb, amsfonts}
\usepackage{geometry}
\usepackage{booktabs}
\geometry{margin=1in}
\title{Haskan2 Rendering Pipeline Mathematical Trace}
\author{Rune}
\date{\today}

\begin{document}
\maketitle

\section{Coordinate Systems}

\subsection{World Space}
\begin{itemize}
\item Right-handed coordinate system
\item +X = right, +Y = geometric up, +Z = forward (toward viewer in default view)
\item \textbf{Visual orientation}: Camera is rotated 180° around its forward axis, so +Y renders at the \textbf{bottom} of the screen and -Y renders at the \textbf{top}. This makes the world appear "upside down" relative to the geometric coordinate system.
\item All procedural meshes (cube, sphere, plane) have outward-facing normals and CCW winding when viewed from outside.
\item glTF models are loaded as-is with standard CCW winding.
\end{itemize}

\subsection{View Space (Custom lookAtNegativeYUp)}
\begin{itemize}
\item \textbf{Right-handed} coordinate system (det = +1)
\item Forward: $-Z$ (points from camera toward scene)
\item Right: $+X$ = world $+X$ (right)
\item Up: $+Y$ = world $+Y$ (up) in the view matrix construction, but the camera is oriented so that visually -Y is up
\item View matrix $\mathbf{V}$ transforms world $\to$ view
\item linear stores matrices in \textbf{row-major} order
\end{itemize}

\textbf{Custom lookAtNegativeYUp} (Camera.hs):
\begin{align}
\mathbf{f} &= \text{normalize}(\mathbf{c} - \mathbf{e}) \quad \text{(forward, points toward scene)} \\
\mathbf{s} &= \text{normalize}(\mathbf{f} \times \mathbf{u}) \quad \text{(right)} \\
\mathbf{u'} &= \mathbf{s} \times \mathbf{f} \quad \text{(true up)}
\end{align}

Where $\mathbf{u} = (0, 1, 0)$ for the orbit math (worldUp), but $\mathbf{u} = (0, -1, 0)$ is passed to lookAtNegativeYUp in orbitalToMatrix to create a 180° rotation around the forward axis.

View matrix (row-major):
\[
\mathbf{V} = \begin{pmatrix}
s_x & s_y & s_z & -\mathbf{s} \cdot \mathbf{e} \\
u'_x & u'_y & u'_z & -\mathbf{u'} \cdot \mathbf{e} \\
-f_x & -f_y & -f_z & \mathbf{f} \cdot \mathbf{e} \\
0 & 0 & 0 & 1
\end{pmatrix}
\]

\textbf{Key point}: The third row is $-\mathbf{f}$ because in the view space, $-Z$ is forward (toward the scene). This is a right-handed convention.

\subsection{The 180° Camera Rotation}
In orbitalToMatrix, we pass $\mathbf{u} = (0, -1, 0)$ to lookAtNegativeYUp instead of $(0, 1, 0)$. This is a 180° rotation around the camera's forward axis ($\mathbf{f}$). Effects:
\begin{itemize}
\item World +Y maps to view -Y → screen \textbf{bottom}
\item World -Y maps to view +Y → screen \textbf{top} (visually "up")
\item World +X maps to view -X → screen \textbf{left}
\item World -X maps to view +X → screen \textbf{right}
\item Winding is preserved (rotation has det = +1, not a reflection)
\item Backface culling works correctly with CCW winding
\end{itemize}

\subsection{Clip Space}
\begin{itemize}
\item After projection: $x, y, z, w$ in homogeneous coordinates
\item Vulkan NDC: $x \in [-1, 1]$ (left to right), $y \in [-1, 1]$ \textbf{top to bottom}
\item After perspective divide: $x_{ndc} = x/w$, $y_{ndc} = y/w$, $z_{ndc} = z/w$
\end{itemize}

\subsection{Framebuffer Space}
\begin{itemize}
\item Origin $(0, 0)$ at \textbf{top-left}
\item $y$ increases \textbf{downward}
\item Viewport transform: $(x_{ndc}, y_{ndc}) \to (x_{screen}, y_{screen})$
\item $x_{screen} = (x_{ndc} + 1) \cdot \frac{width}{2}$
\item $y_{screen} = (y_{ndc} + 1) \cdot \frac{height}{2}$ (so $y_{ndc} = -1 \to y_{screen} = 0$ = top)
\end{itemize}

\section{Vertex Transformation Pipeline}

\subsection{Step 1: Model to World}
Model matrix $\mathbf{M}$ is identity for most objects.
\[
\mathbf{p}_{world} = \mathbf{M} \cdot \mathbf{p}_{model} = \mathbf{p}_{model}
\]

\subsection{Step 2: World to View}
View matrix $\mathbf{V}$ from lookAtNegativeYUp:
\[
\mathbf{p}_{view} = \mathbf{V} \cdot \mathbf{p}_{world}
\]

The custom lookAtNegativeYUp builds a right-handed view matrix (det = +1). When called with $\mathbf{u} = (0, -1, 0)$, it produces a view matrix that is rotated 180° around the forward axis compared to standard lookAt.

\subsection{Step 3: View to Clip — Standard Projection}
Projection matrix $\mathbf{P}$ uses standard perspective projection \textbf{without} Y-flip:
\[
\mathbf{P} = \mathbf{P}_{persp}
\]

The Y-down behavior is achieved via \textbf{negative viewport height} (Vulkan 1.1+ core):

\begin{verbatim}
VkViewport viewport = {
    .x = 0,
    .y = height,           // Start from bottom
    .width = width,
    .height = -height,     // Negative height flips Y
    .minDepth = 0.0,
    .maxDepth = 1.0
};
\end{verbatim}

This achieves the same Y-down result without modifying the projection matrix or reversing winding. Benefits:
\begin{itemize}
\item Standard projection matrix (no Y-flip)
\item CCW meshes remain CCW in clip space
\item Backface culling works with $\text{frontFace} = \text{CCW}$
\item Compatible with all external assets (glTF, OBJ, etc.)
\end{itemize}

\subsection{Current State (Verified)}
\begin{itemize}
\item Projection matrix: Standard perspective (no Y-flip)
\item Viewport: Negative height ($\text{height} = -h$, $y = h$)
\item Culling: $\text{VK\_CULL\_MODE\_BACK\_BIT}$ enabled
\item Front face: $\text{VK\_FRONT\_FACE\_COUNTER\_CLOCKWISE}$ (standard)
\item All procedural meshes have outward normals with CCW winding
\item glTF models load with standard CCW winding and render correctly
\end{itemize}

\subsection{Step 4: Perspective Divide}
\[
\mathbf{p}_{ndc} = \begin{pmatrix} x_{clip}/w_{clip} \\ y_{clip}/w_{clip} \\ z_{clip}/w_{clip} \end{pmatrix}
\]

\subsection{Step 5: Viewport Transform}
\[
x_{screen} = \frac{x_{ndc} + 1}{2} \cdot width, \quad
y_{screen} = \frac{y_{ndc} + 1}{2} \cdot height
\]

\section{Historical Note: The Y-Flip and Winding Problem}

\subsection{The Problem (Resolved)}
Earlier versions used a Y-flip in the projection matrix:
\[
\mathbf{P} = \mathbf{Y}_{flip} \cdot \mathbf{P}_{persp}, \quad
\mathbf{Y}_{flip} = \text{diag}(1, -1, 1, 1)
\]

This is a reflection (det = -1) that reversed triangle winding in clip space:
\begin{itemize}
\item CCW in view space $\to$ CW in clip space
\item Broke backface culling for standard CCW meshes
\end{itemize}

\subsection{The Wrong Fix}
We temporarily used $\text{frontFace} = \text{VK\_FRONT\_FACE\_CLOCKWISE}$ and disabled culling ($\text{VK\_CULL\_MODE\_NONE}$) as workarounds. These were wrong because they broke glTF compatibility and papered over the real problem.

\subsection{The Correct Fix}
Removed Y-flip from projection. Use negative viewport height for Vulkan Y-down. Enable standard backface culling with CCW front face. All meshes (procedural and glTF) now render correctly.

\section{Camera Orbital Math}

\subsection{worldUp Definition}
\begin{verbatim}
worldUp :: V3 CFloat
worldUp = V3 0 1 0
\end{verbatim}

This is the geometric up vector used for:
\begin{itemize}
\item Azimuth (yaw) rotation axis
\item Right vector computation: $\mathbf{r} = \text{normalize}(\text{worldUp} \times \mathbf{fwd})$
\item MoveRight direction
\end{itemize}

\subsection{Azimuth and Elevation}
\begin{itemize}
\item Azimuth: Rotation around worldUp (yaw). Positive = counter-clockwise when viewed from above.
\item Elevation: Rotation around local right axis (pitch). Positive = looking \textbf{up} (toward world +Y).
\item Default camera: azimuth = 0°, elevation = 30° (looking 30° above horizon)
\end{itemize}

\subsection{Camera Orientation}
With the 180° rotation in lookAtNegativeYUp:
\begin{itemize}
\item Default camera at elevation = 30° looks \textbf{downward} toward the scene
\item World +Y (green) appears at screen bottom
\item World -Y (yellow) appears at screen top
\item Mouse Y axis is inverted to match this orientation
\end{itemize}

\section{Skybox Ray Computation}

\subsection{Current Implementation}
\begin{verbatim}
computeSkyboxRays view proj =
  let V4 (V4 v00 v01 v02 _) (V4 v10 v11 v12 _) (V4 v20 v21 v22 _) _ = view
      worldRot = V3 (V3 v00 v01 v02) (V3 v10 v11 v12) (V3 v20 v21 v22)
      
      V4 (V4 fx _ _ _) (V4 _ fy _ _) _ _ = proj
      
      ndcToDir (V4 x y _z _w) =
        let viewDir = V3 (x / fx) (y / fy) (-1)
        in normalize (worldRot !* viewDir)
\end{verbatim}

\subsection{Analysis of worldRot}
Since linear stores matrices in row-major, the view matrix rows are:
\begin{itemize}
\item Row 0: $(v_{00}, v_{01}, v_{02}) = \mathbf{s}$ (world-space right vector)
\item Row 1: $(v_{10}, v_{11}, v_{12}) = \mathbf{u'}$ (world-space up vector)
\item Row 2: $(v_{20}, v_{21}, v_{22}) = -\mathbf{f}$ (world-space forward vector)
\end{itemize}

So:
\[
\mathbf{worldRot} = \begin{pmatrix}
| & | & | \\
\mathbf{s} & \mathbf{u'} & -\mathbf{f} \\
| & | & |
\end{pmatrix}
\]

This transforms view-space directions to world-space directions.

\textbf{Correction}: Earlier versions extracted \textbf{columns} instead of rows, which gave $\mathbf{R}$ (world$\to$view) instead of $\mathbf{R}^T$ (view$\to$world). This caused a severe diagonal tilt for non-axis-aligned cameras.

\subsection{Analysis of ndcToDir}
For a pixel at NDC $(x, y)$:
\[
\mathbf{v}_{view} = \begin{pmatrix} x/f_x \\ y/f_y \\ -1 \end{pmatrix}
\]

Then:
\[
\mathbf{v}_{world} = \text{normalize}(\mathbf{worldRot} \cdot \mathbf{v}_{view})
\]

\textbf{Why no $-x$ negation?}
Because lookAtNegativeYUp produces a \textbf{right-handed} view space where $+X$ points to world $+X$ (right). Screen-right ($+x$) should look toward world-right ($+\mathbf{s}$), so we use $+x/f_x$.

\textbf{Why $-1$ for Z?}
In view space, $-Z$ points toward the scene (forward). We want the ray to point forward, so we use $-1$.

\subsection{Verification}
At screen center $(x, y) = (0, 0)$:
\[
\mathbf{v}_{view} = (0, 0, -1)
\]
\[
\mathbf{v}_{world} = \text{normalize}(-\mathbf{f}) = \mathbf{f}
\]
This is correct: center of screen looks in the forward direction.

At screen right $(x, y) = (1, 0)$:
\[
\mathbf{v}_{view} = (1/f_x, 0, -1)
\]
\[
\mathbf{v}_{world} = \text{normalize}(\frac{1}{f_x}\mathbf{s} - \mathbf{f})
\]

Since $\mathbf{s}$ points to world-right, screen-right looks toward world-right. This is correct.

\section{Mesh Conventions}

\subsection{Procedural Meshes}
All procedural meshes use \textbf{outward-facing normals} and \textbf{CCW winding} when viewed from outside:

\begin{tabular}{lll}
\toprule
Mesh & Normal Direction & Winding \\
\midrule
unitCube & Outward on all 6 faces & CCW from outside \\
uvSphere & Radially outward & CCW from outside \\
uvPlane & $(0, -1, 0)$ (facing down) & CCW from below \\
groundPlane & $(0, -1, 0)$ (facing down) & CCW from below \\
\bottomrule
\end{tabular}

\subsection{UV Mapping}
\begin{itemize}
\item Vulkan convention: $(0, 0)$ = top-left, $(1, 1)$ = bottom-right
\item unitCube: UVs flipped horizontally to match visual orientation
\item uvSphere: V flipped ($V = 1 - \text{latIdx}/\text{latCount}$) so poles map correctly
\item glTF: UVs pass through unchanged (Vulkan convention matches glTF)
\end{itemize}

\section{Cubemap Face Convention}

\subsection{Vulkan Cubemap Faces}
Faces are indexed 0--5:
\begin{enumerate}
\item[0] +X (right)
\item[1] -X (left)
\item[2] +Y (up)
\item[3] -Y (down)
\item[4] +Z (forward)
\item[5] -Z (backward)
\end{enumerate}

\subsection{Face Orientation (Vulkan Spec)}
For face +X (viewed from +X looking toward origin):
\begin{itemize}
\item Image $+U$ direction = $-Z$ (backward)
\item Image $+V$ direction = $-Y$ (down)
\end{itemize}

For face -X (viewed from -X looking toward origin):
\begin{itemize}
\item Image $+U$ direction = $+Z$ (forward)
\item Image $+V$ direction = $-Y$ (down)
\end{itemize}

For face +Y (viewed from +Y looking toward origin):
\begin{itemize}
\item Image $+U$ direction = $+X$ (right)
\item Image $+V$ direction = $+Z$ (forward)
\end{itemize}

For face -Y (viewed from -Y looking toward origin):
\begin{itemize}
\item Image $+U$ direction = $+X$ (right)
\item Image $+V$ direction = $-Z$ (backward)
\end{itemize}

For face +Z (viewed from +Z looking toward origin):
\begin{itemize}
\item Image $+U$ direction = $+X$ (right)
\item Image $+V$ direction = $-Y$ (down)
\end{itemize}

For face -Z (viewed from -Z looking toward origin):
\begin{itemize}
\item Image $+U$ direction = $-X$ (left)
\item Image $+V$ direction = $-Y$ (down)
\end{itemize}

\subsection{Cubemap Directory Structure}
Cubemaps are loaded from \texttt{data/textures/cubemaps/\{name\}/}:
\begin{verbatim}
data/textures/cubemaps/
├── debug/           # Colored test cubemap
│   ├── posx.png     # Red (+X)
│   ├── negx.png     # Blue (-X)
│   ├── posy.png     # Green (+Y)
│   ├── negy.png     # Yellow (-Y)
│   ├── posz.png     # Gray (+Z)
│   └── negz.png     # Brown (-Z)
└── env1/            # Real environment map
    ├── posx.png
    ├── negx.png
    ├── posy.png
    ├── negy.png
    ├── posz.png
    └── negz.png
\end{verbatim}

\subsection{cmft Output Convention}
cmft generates faces in the \textbf{OpenGL} convention by default. The OpenGL convention differs from Vulkan:

\textbf{OpenGL vs Vulkan differences:}
\begin{itemize}
\item For +X and -X faces: OpenGL has $+U$ pointing in the \textbf{opposite} Z direction compared to Vulkan
\item This means +X and -X faces need to be \textbf{horizontally flipped} when loading into Vulkan
\item For +Y face: OpenGL has $+V$ pointing in a different direction than Vulkan
\end{itemize}

\textbf{Fixed}: Real environment cubemap faces (env1) were processed to match Vulkan convention by flipping vertically and swapping/mirroring all axis pairs.

\section{Resolved Issues}

\subsection{Issue 1: Cubemap Face Orientation Mismatch}
\textbf{Status}: Resolved. env1 cubemap faces were flipped and swapped to match Vulkan convention.

\subsection{Issue 2: Screenshot Y-Flip}
\textbf{Status}: Resolved. Screenshot capture now handles the negative viewport correctly.

\section{Critical Invariants}

\begin{enumerate}
\item View matrix must have det = +1 (right-handed, no reflection)
\item All meshes (procedural and imported) use CCW winding with outward normals
\item Projection matrix is standard (no Y-flip)
\item Viewport uses negative height for Vulkan Y-down
\item Backface culling enabled with CCW front face
\item Camera up vector for orbit: worldUp = V3 0 1 0
\item Camera up vector for lookAt: V3 0 (-1) 0 (180° rotation around forward)
\end{enumerate}

\end{document}